\section*{Dag 4\&5. 08-01-2026/09-01-2026}
\vspace{1em}

\textbf{Exercise 5.1 - Detecting a Joystick Interaction}
\begin{lstlisting}
#include "stm32f30x_conf.h"
#include "30010_io.h"
#include <stdint.h>

// --------------------------------------------------
// Read joystick state
// --------------------------------------------------
uint8_t readJoystick(void)
{
    uint8_t state = 0;

    if (GPIOA->IDR & (1 << 4)) state |= (1 << 0); // UP
    if (GPIOB->IDR & (1 << 0)) state |= (1 << 1); // DOWN
    if (GPIOC->IDR & (1 << 1)) state |= (1 << 2); // LEFT
    if (GPIOC->IDR & (1 << 0)) state |= (1 << 3); // RIGHT
    if (GPIOB->IDR & (1 << 5)) state |= (1 << 4); // CENTER

    return state;
}

// --------------------------------------------------
// Simple delay
// --------------------------------------------------
void delay_ms(uint32_t ms)
{
    for (uint32_t i = 0; i < ms * 8000; i++) {
        __NOP();
    }
}

// --------------------------------------------------
// Main
// --------------------------------------------------
int main(void)
{
    uart_init(9600);

    // Enable GPIO clocks
    RCC->AHBENR |= RCC_AHBPeriph_GPIOA
                | RCC_AHBPeriph_GPIOB
                | RCC_AHBPeriph_GPIOC;

    // Configure joystick pins as INPUT, NO pull-up/down

    // PA4 - UP
    GPIOA->MODER &= ~(0x3 << (4 * 2));
    GPIOA->PUPDR &= ~(0x3 << (4 * 2));

    // PB0 - DOWN
    GPIOB->MODER &= ~(0x3 << (0 * 2));
    GPIOB->PUPDR &= ~(0x3 << (0 * 2));

    // PB5 - CENTER
    GPIOB->MODER &= ~(0x3 << (5 * 2));
    GPIOB->PUPDR &= ~(0x3 << (5 * 2));

    // PC0 - RIGHT
    GPIOC->MODER &= ~(0x3 << (0 * 2));
    GPIOC->PUPDR &= ~(0x3 << (0 * 2));

    // PC1 - LEFT
    GPIOC->MODER &= ~(0x3 << (1 * 2));
    GPIOC->PUPDR &= ~(0x3 << (1 * 2));

    uint8_t lastState = 0;

    while (1)
    {
        uint8_t state = readJoystick();

        if (state != lastState)
        {
            if (state == 0) {
                printf("RELEASED\r\n");
            } else {
                if (state & (1 << 0)) printf("UP\r\n");
                if (state & (1 << 1)) printf("DOWN\r\n");
                if (state & (1 << 2)) printf("LEFT\r\n");
                if (state & (1 << 3)) printf("RIGHT\r\n");
                if (state & (1 << 4)) printf("CENTER\r\n");
            }

            lastState = state;
            delay_ms(200);   // simple debounce / lockout
        }
    }
}
\end{lstlisting}

\textbf{Exercise 5.2 - Controlling the RGB LED}
\begin{lstlisting}
    #include "stm32f30x_conf.h"
#include "30010_io.h"
#include <stdint.h>

// LED color definitions
#define LED_OFF     0
#define LED_RED     1
#define LED_GREEN   2
#define LED_BLUE    3
#define LED_YELLOW  4
#define LED_MAGENTA 5
#define LED_CYAN    6
#define LED_WHITE   7

uint8_t readJoystick(void)
{
    uint8_t state = 0;

    if (GPIOA->IDR & (1 << 4)) state |= (1 << 0); // UP
    if (GPIOB->IDR & (1 << 0)) state |= (1 << 1); // DOWN
    if (GPIOC->IDR & (1 << 1)) state |= (1 << 2); // LEFT
    if (GPIOC->IDR & (1 << 0)) state |= (1 << 3); // RIGHT
    if (GPIOB->IDR & (1 << 5)) state |= (1 << 4); // CENTER

    return state;
}

void setLed(uint8_t color)
{
    // ACTIVE-LOW LED: 1 = OFF, 0 = ON
    GPIOB->ODR |= (1 << 4); // Red off
    GPIOC->ODR |= (1 << 7); // Green off
    GPIOA->ODR |= (1 << 9); // Blue off

    switch (color)
    {
        case LED_RED:
            GPIOB->ODR &= ~(1 << 4);
            break;

        case LED_GREEN:
            GPIOC->ODR &= ~(1 << 7);
            break;

        case LED_BLUE:
            GPIOA->ODR &= ~(1 << 9);
            break;

        case LED_YELLOW:
            GPIOB->ODR &= ~(1 << 4);
            GPIOC->ODR &= ~(1 << 7);
            break;

        case LED_MAGENTA:
            GPIOB->ODR &= ~(1 << 4);
            GPIOA->ODR &= ~(1 << 9);
            break;

        case LED_CYAN:
            GPIOC->ODR &= ~(1 << 7);
            GPIOA->ODR &= ~(1 << 9);
            break;

        case LED_WHITE:
            GPIOB->ODR &= ~(1 << 4);
            GPIOC->ODR &= ~(1 << 7);
            GPIOA->ODR &= ~(1 << 9);
            break;

        default:
            break;
    }
}

int main(void)
{
    uart_init(9600);

    RCC->AHBENR |= RCC_AHBPeriph_GPIOA
                | RCC_AHBPeriph_GPIOB
                | RCC_AHBPeriph_GPIOC;

    GPIOA->MODER &= ~(0x3 << (4 * 2));
    GPIOB->MODER &= ~(0x3 << (0 * 2));
    GPIOB->MODER &= ~(0x3 << (5 * 2));
    GPIOC->MODER &= ~(0x3 << (0 * 2));
    GPIOC->MODER &= ~(0x3 << (1 * 2));

    GPIOB->MODER &= ~(0x3 << (4 * 2));
    GPIOB->MODER |=  (0x1 << (4 * 2));

    GPIOC->MODER &= ~(0x3 << (7 * 2));
    GPIOC->MODER |=  (0x1 << (7 * 2));

    GPIOA->MODER &= ~(0x3 << (9 * 2));
    GPIOA->MODER |=  (0x1 << (9 * 2));

    uint8_t lastState = 0;

    while (1)
    {
        uint8_t state = readJoystick();

        if (state != lastState)
        {
            if (state & (1 << 0))       setLed(LED_RED);     // UP
            else if (state & (1 << 1))  setLed(LED_BLUE);    // DOWN
            else if (state & (1 << 2))  setLed(LED_GREEN);   // LEFT
            else if (state & (1 << 3))  setLed(LED_YELLOW);  // RIGHT
            else if (state & (1 << 4))  setLed(LED_WHITE);   // CENTER
            else                        setLed(LED_OFF);

            lastState = state;
        }
    }
}
\end{lstlisting}

\textbf{Exercise 5.3: Translate the following C code to phyton}

Translate the following C code to Python. C code:
\begin{lstlisting}
    struct node
        {
        int i;
        int j;
        };
    int main(void)
        {
        struct node s;
        s.i = 10;
        s.j = 5;
        printf ("%i",s.i);
        printf ("\n%i",s.j);
        }
\end{lstlisting}

Translated to Python:
\begin{lstlisting}
    class node:
        def __init__(self,i,j):
            self.i = i
            self.j = j s = node(10,5)
            print("%d" % s.i)
            print("%d" % s.j
\end{lstlisting}

\textbf{Exercise 6.1 – Timers}


